\subsection{ResultBoundaryFluxes.txt}
%
\subsubsection{Purpose}
This file stores the instantaneous discharge crossing each boundary every \param{dt\_time\_series}. In case of the boundary is composed of several edges, the stored discharge is the sum of the discharge crossing each edge. As the unit vector normal to the boundary edges is oriented from the exterior to the interior, a positive discharge enters the domain.
%
\subsubsection{Format description}
\texttt{<Tab>}: tabulation character. \\
    \begin{tabular}{p{6cm} p{9cm}}
    \hline
    \texttt{time <Tab> BC\_$x_1$ <Tab> \ldots <Tab> BC\_$x_N$} & This is the header line. $x_1$, \ldots, $x_N$ represent the boundary numbers as defined in the \file{2dm} file \\
    \hline 
    $t$ \texttt{<Tab>} $Q_1$ \texttt{<Tab>} \ldots \texttt{<Tab>} $Q_N$ & $t$ is the stored time, $Q_i$ is the discharge (in $\mathrm{m^3.s^{-1}}$) through the boundary $x_i$ \\
    \hline
    \ldots & one line per stored time \\ \hline
    \end{tabular}
%
\subsubsection{Example}
\begin{verbatim}
time	BC_1	BC_2	BC_3
2.0000000000e+1	0.0000000000e+0	4.3274586317e+1	-1.4193271476e+1
3.0000000000e+1	0.0000000000e+0	4.3149535653e+1	-1.4227591312e+1
4.0000000000e+1	0.0000000000e+0	4.3074013841e+1	-1.4246148173e+1
5.0000000000e+1	0.0000000000e+0	4.3030038602e+1	-1.4258119344e+1
6.0000000000e+1	0.0000000000e+0	4.3000513361e+1	-1.4266604580e+1
\end{verbatim}